We describe a Lean 4 formalization of two finite-injury priority arguments
over a single, shared oracle-machine foundation: the Friedberg--Muchnik
theorem, which solves Post's problem by constructing two mutually
nonreducible computably enumerable sets, and the Sacks Splitting Theorem in
its splitting-with-non-reducibility form, which splits any non-computable
c.e.\ set into two c.e.\ halves computing neither.

The development is \emph{analytic} rather than synthetic: oracle
computation is not a black-box semantic relation but an explicit language
of oracle programs equipped with a step-indexed evaluator over finite
oracle information.  A halting bounded evaluation records both its output
and a strict upper bound on the oracle positions it queried, and the
central use theorem proves that such a computation is preserved under extra
fuel and under oracle agreement below the recorded use.  Restraint,
injury, initialization, attention, eventual stability and effective stage
construction are all represented directly, and computability of each stage
function is discharged against Mathlib's ordinary primitive- and
partial-recursive infrastructure.

The second theorem is a test of the first's foundation.  The whole
foundation transferred unchanged, and the recorded use---a design decision
taken for Friedberg--Muchnik---turned out to be what bounds the restraint
in Sacks, a theorem it was not designed for.  The priority layer did not
transfer at all: Friedberg--Muchnik's injury argument counts discrete
events, and a splitting requirement never acts, so a different induction is
required.  We report this as a negative result: ``finite injury'' names a
family of arguments rather than a single reusable lemma.  The two libraries
are siblings, importing a common foundation and not each other, and the
build enforces it.

Finally, we relate the local reducibility relation to Mathlib's own
\texttt{RecursiveIn} model by proving that every reduction there is
realized by a program of our language.  Since both conclusions are
negations of reducibility, this is the direction that rules out a
vacuously weak relation, and both theorems are restated over Mathlib's
notion.  The development builds with no \texttt{sorry}, and every headline
theorem depends only on \texttt{propext}, \texttt{Classical.choice} and
\texttt{Quot.sound}.
